% Fichier engendré par tools/generate-tex.py à partir de 10-specialite-mathematiques-terminale-s.
% Les démonstrations en français sont transcrites du fichier Lean (skill transcribe-lean-proof).
\documentclass[11pt,a4paper]{article}

% Compile aussi bien avec pdflatex qu'avec xelatex ou tectonic.
\usepackage{iftex}
\ifPDFTeX
  \usepackage[T1]{fontenc}
  \usepackage[utf8]{inputenc}
\else
  \usepackage{fontspec}
\fi
\usepackage[french]{babel}
\usepackage{amsmath,amssymb,amsthm}
\usepackage[margin=2.5cm]{geometry}
\usepackage[hidelinks]{hyperref}

% Les figures sont dessinées en TikZ : elles reprennent ainsi les
% polices du document. Voir la skill `illustrate-theorem`.
\usepackage{tikz}
\usetikzlibrary{calc}

\theoremstyle{definition}
\newtheorem{definition}{Définition}
\newtheorem{exemple}{Exemple}
\theoremstyle{plain}
\newtheorem{theoreme}{Théorème}
\newtheorem{lemme}[theoreme]{Lemme}

% Un exemple ne se démontre pas : on explique ce qu'il montre.
\newenvironment{explication}{\begin{proof}[Explication]}{\end{proof}}

% Un énoncé écrit mais non démontré : la preuve Lean est un `sorry`.
\newcommand{\admis}{\par\smallskip\noindent\textit{Admis : l'énoncé est écrit en Lean, sa démonstration reste à faire.}\par}

% Renvoi à la déclaration Lean, une ligne après la démonstration, aligné à droite.
\newcommand{\source}[2]{\par\smallskip\noindent\hfill{\footnotesize\href{#1}{\texttt{#2}}}\par\smallskip}

\title{Arithmétique, matrices et graphes}
\date{}

\begin{document}
\maketitle
% Les identifiants Lean en machine à écrire ne se coupent pas : on tolère des
% espacements plus lâches plutôt que des lignes qui débordent.
\sloppy

\section{Arithmetique}

\noindent
Spécialité de terminale S, section \og{} Arithmétique \fg{}. Les entiers relatifs sont
ceux de Mathlib : $a \mid b$ se lit \og{} $a$ divise $b$ \fg{}, $a \bmod b$ est le reste de
la division euclidienne, et $a \equiv b \ [n]$ la congruence modulo $n$.

\medskip\noindent
C'est le chapitre du programme le plus proche de Mathlib : presque tous ses énoncés y
existent déjà, parfois sous un autre nom. L'exercice n'est donc plus de démontrer mais de
retrouver la formulation exacte et de vérifier qu'elle dit bien la même chose que l'énoncé
français \textemdash{} exercice différent, et pas toujours plus facile.

\subsection{Division euclidienne}

\begin{theoreme}
Division euclidienne dans $\mathbb{Z}$ : pour $b > 0$, tout entier $a$ s'écrit d'une seule
façon
\[ a = bq + r , \qquad 0 \le r < b . \]
\end{theoreme}

\begin{proof}
\emph{Existence.} Le couple $(a \div b,\ a \bmod b)$ convient : la bibliothèque démontre
l'identité $a \bmod b + b\,(a \div b) = a$, ainsi que $0 \le a \bmod b$ et
$a \bmod b < b$ pour $b > 0$.

\emph{Unicité.} Soit $(q, r)$ un autre couple convenable. La bibliothèque donne
précisément la caractérisation cherchée : pour $b \ne 0$,
\[ a \div b = q \ \text{ et } \ a \bmod b = r
   \iff r + bq = a \ \text{ et } \ 0 \le r < |b| . \]
Comme $b > 0$, sa valeur absolue est $b$ lui-même, et les trois conditions dont on dispose
sont exactement le membre de droite. Le couple $(q, r)$ est donc $(a \div b,\ a \bmod b)$.
\end{proof}
\source{https://github.com/Commutator-IO/learn-lean/blob/main/courses/02-lycee/10-specialite-mathematiques-terminale-s/Arithmetique.lean\#L16}{Arithmetique.lean\#L16}

\subsection{Congruences}

\begin{theoreme}
Les congruences sont compatibles avec la somme et le produit :
\[ a \equiv b \ [n] \ \text{ et } \ c \equiv d \ [n]
   \implies a + c \equiv b + d \ [n] \ \text{ et } \ ac \equiv bd \ [n] . \]
\end{theoreme}

\begin{proof}
Dire que $a \equiv b \ [n]$, c'est dire que $n$ divise $b - a$. Les deux règles sont alors
des calculs de divisibilité.

Pour la somme :
\[ (b + d) - (a + c) = (b - a) + (d - c) , \]
somme de deux multiples de $n$, donc multiple de $n$.

Pour le produit, on fait apparaître les deux différences en insérant un terme :
\[ bd - ac = b\,(d - c) + c\,(b - a) , \]
et chacun des deux termes est un multiple de $n$.

C'est toute la raison d'être du calcul modulaire : les congruences sont compatibles avec la
structure d'anneau de $\mathbb{Z}$, donc on peut remplacer un nombre par son reste avant de
calculer plutôt qu'après.
\end{proof}
\source{https://github.com/Commutator-IO/learn-lean/blob/main/courses/02-lycee/10-specialite-mathematiques-terminale-s/Arithmetique.lean\#L31}{Arithmetique.lean\#L31}

\begin{theoreme}
Critère de divisibilité par $9$ : un entier est congru à la somme de ses chiffres
modulo $9$, et il est divisible par $9$ si et seulement si cette somme l'est.
\end{theoreme}

\begin{proof}
Le premier point est le cas $b = 9$ d'un énoncé plus large : si la base $b'$ vérifie
$b' \equiv 1$ modulo $b$, alors tout entier est congru modulo $b$ à la somme de ses chiffres
en base $b'$. Ici $10 \equiv 1 \ [9]$, vérification numérique immédiate.

L'argument est le suivant : un nombre écrit $\overline{c_k \dots c_0}$ vaut
$\sum_i c_i 10^{i}$, et chaque $10^{i}$ est congru à $1^{i} = 1$ modulo $9$, de sorte que
la somme est congrue à $\sum_i c_i$.

Le critère de divisibilité s'en déduit : un nombre est divisible par $9$ exactement quand
son reste modulo $9$ est nul, donc quand celui de la somme de ses chiffres l'est.
\end{proof}
\source{https://github.com/Commutator-IO/learn-lean/blob/main/courses/02-lycee/10-specialite-mathematiques-terminale-s/Arithmetique.lean\#L37}{Arithmetique.lean\#L37}

\subsection{PGCD et algorithme d'Euclide}

\begin{theoreme}
Algorithme d'Euclide :
\[ \mathrm{pgcd}(a, b) = \mathrm{pgcd}(b \bmod a,\ a) . \]
\end{theoreme}

\begin{proof}
Pour $a = 0$, les deux membres valent $b$ : le pgcd de $0$ et $b$ est $b$, et
$b \bmod 0 = b$.

Pour $a > 0$, c'est l'équation de récurrence qui définit le pgcd dans la bibliothèque. Elle
fait décroître strictement le premier argument, puisque $b \bmod a < a$ : c'est ce qui
assure la terminaison de l'algorithme, et c'est aussi ce qui justifie que la définition
soit bien posée.
\end{proof}
\source{https://github.com/Commutator-IO/learn-lean/blob/main/courses/02-lycee/10-specialite-mathematiques-terminale-s/Arithmetique.lean\#L46}{Arithmetique.lean\#L46}

\subsection{Théorèmes de Bézout et de Gauss}

\begin{theoreme}
Théorème de Bézout :
\[ \mathrm{pgcd}(a, b) = 1 \iff \exists (u, v) \in \mathbb{Z}^{2},\ ua + vb = 1 . \]
\end{theoreme}

\begin{proof}
La bibliothèque appelle \emph{copremiers} deux éléments $a$ et $b$ pour lesquels il existe
$u$ et $v$ avec $ua + vb = 1$, et démontre que, dans $\mathbb{Z}$, cette propriété équivaut à
$\mathrm{pgcd}(a, b) = 1$. Les deux sens ne coûtent pas le même prix.

\emph{Une combinaison force le pgcd à valoir un.} Si $ua + vb = 1$, alors tout diviseur
commun de $a$ et $b$ divise le membre de gauche, donc divise $1$ : le pgcd vaut $1$. C'est
immédiat.

\emph{Le pgcd fournit une combinaison.} C'est le sens qui a du contenu, et la bibliothèque
ne l'obtient pas par un argument d'existence : elle \emph{calcule} les coefficients, par
l'algorithme d'Euclide étendu. À chaque division $b = aq + r$, on remonte une écriture de
$r$ comme combinaison de $a$ et $b$ ; en fin de course, le dernier reste non nul est le
pgcd, et l'on dispose de son écriture. L'égalité $\mathrm{pgcd}(a,b) = au + bv$ vaut donc
toujours, et l'hypothèse $\mathrm{pgcd} = 1$ ne fait que la spécialiser.

Une fois l'équivalence en main, les deux membres de l'énoncé sont le même, et chaque sens se
démontre en reprenant le couple $(u, v)$ fourni par l'autre.
\end{proof}
\source{https://github.com/Commutator-IO/learn-lean/blob/main/courses/02-lycee/10-specialite-mathematiques-terminale-s/Arithmetique.lean\#L55}{Arithmetique.lean\#L55}

\begin{theoreme}
Théorème de Gauss : si $a \mid bc$ et $\mathrm{pgcd}(a, b) = 1$, alors $a \mid c$.
\end{theoreme}

\begin{proof}
L'argument est Bézout : de $ua + vb = 1$ on tire $uac + vbc = c$, et $a$ divise
les deux termes de gauche \textemdash{} le premier visiblement, le second parce que
$a \mid bc$.

\medskip\noindent
L'énoncé formel est écrit pour des éléments \emph{copremiers} ; on y accède en traduisant
l'hypothèse $\mathrm{pgcd}(a, b) = 1$ en coprimalité, par l'équivalence du théorème
précédent.
\end{proof}
\source{https://github.com/Commutator-IO/learn-lean/blob/main/courses/02-lycee/10-specialite-mathematiques-terminale-s/Arithmetique.lean\#L62}{Arithmetique.lean\#L62}

\subsection{Équations diophantiennes}

\begin{theoreme}
L'équation $ax + by = c$ a une solution entière si et seulement si
$\mathrm{pgcd}(a, b) \mid c$.
\end{theoreme}

\begin{proof}
\emph{Sens direct.} Si $ax + by = c$, alors $\mathrm{pgcd}(a,b)$ divise $a$ et $b$, donc
$ax$ et $by$, donc leur somme $c$.

\emph{Sens réciproque.} Écrivons $c = \mathrm{pgcd}(a, b) \times k$. L'identité de Bézout
de la bibliothèque fournit deux coefficients $u$ et $v$ tels que
\[ \mathrm{pgcd}(a, b) = au + bv , \]
et en multipliant par $k$ :
\[ a\,(uk) + b\,(vk) = \mathrm{pgcd}(a, b)\,k = c . \]
Le couple $(uk, vk)$ convient.
\end{proof}
\source{https://github.com/Commutator-IO/learn-lean/blob/main/courses/02-lycee/10-specialite-mathematiques-terminale-s/Arithmetique.lean\#L68}{Arithmetique.lean\#L68}

\begin{theoreme}
Forme des solutions dans le cas $\mathrm{pgcd}(a, b) = 1$ et $b \ne 0$ : si $(x_0, y_0)$
est une solution particulière de $ax + by = c$, toute solution $(x, y)$ s'écrit
\[ x = x_0 + kb , \qquad y = y_0 - ka , \qquad k \in \mathbb{Z} . \]
\end{theoreme}

\begin{proof}
En retranchant les deux équations $ax + by = c$ et $ax_0 + by_0 = c$ :
\[ b\,(y_0 - y) = a\,(x - x_0) . \]
Ainsi $b$ divise $a(x - x_0)$ ; comme $b$ et $a$ sont premiers entre eux, le théorème de
Gauss donne $b \mid (x - x_0)$. On écrit donc $x - x_0 = kb$, c'est-à-dire
$x = x_0 + kb$.

En reportant dans l'égalité précédente :
\[ b\,(y_0 - y) = a\,kb = b\,(ka) , \]
et comme $b \ne 0$, on peut simplifier : $y_0 - y = ka$, soit $y = y_0 - ka$.

\medskip\noindent
L'hypothèse $b \ne 0$ n'est pas cosmétique. Pour $b = 0$ l'équation devient $ax = c$, dont
les solutions ont un $y$ absolument quelconque, et la description ci-dessus, qui fait
dépendre $x$ et $y$ d'un même paramètre $k$, ne rend plus compte de cette liberté.
\end{proof}
\source{https://github.com/Commutator-IO/learn-lean/blob/main/courses/02-lycee/10-specialite-mathematiques-terminale-s/Arithmetique.lean\#L81}{Arithmetique.lean\#L81}

\subsection{Nombres premiers}

\begin{theoreme}
Il existe une infinité de nombres premiers : au-delà de tout entier $n$, il y en a encore
un.
\end{theoreme}

\begin{proof}
L'argument est celui d'Euclide : un facteur premier de
$n! + 1$ est nécessairement plus grand que $n$, puisque aucun entier compris entre $2$ et
$n$ ne divise $n! + 1$ \textemdash{} chacun divise $n!$ et laisserait un reste $1$.
\end{proof}
\source{https://github.com/Commutator-IO/learn-lean/blob/main/courses/02-lycee/10-specialite-mathematiques-terminale-s/Arithmetique.lean\#L99}{Arithmetique.lean\#L99}

\begin{theoreme}
Décomposition en facteurs premiers : pour $n \ne 0$, le produit de la liste des facteurs
premiers de $n$ vaut $n$, et toute liste de nombres premiers de produit $n$ est une
permutation de celle-là.
\end{theoreme}

\begin{proof}
La liste des facteurs premiers a bien $n$ pour produit, et l'unicité s'énonce comme une égalité de listes \emph{à l'ordre près},
ce qui est la formulation exacte du \og{} unique à l'ordre des facteurs près \fg{} du
programme.

\medskip\noindent
Le programme du lycée \emph{admet} ce théorème \textemdash{} et le collège aussi, où il
figure déjà. Ici il est démontré, ou plutôt emprunté : la démonstration est dans Mathlib,
et repose sur le lemme d'Euclide, lui-même conséquence du théorème de Gauss.
\end{proof}
\source{https://github.com/Commutator-IO/learn-lean/blob/main/courses/02-lycee/10-specialite-mathematiques-terminale-s/Arithmetique.lean\#L104}{Arithmetique.lean\#L104}

\subsection{Petit théorème de Fermat}

\begin{theoreme}
Petit théorème de Fermat : si $p$ est premier et si $a$ n'est pas nul modulo $p$, alors
\[ a^{\,p-1} = 1 \quad \text{dans } \mathbb{Z}/p\mathbb{Z} . \]
\end{theoreme}

\begin{proof}
L'énoncé formel est écrit dans l'anneau $\mathbb{Z}/p\mathbb{Z}$ plutôt qu'avec des congruences d'entiers : \og{} $p$ ne divise
pas $a$ \fg{} devient \og{} $a \ne 0$ \fg{}, et \og{} $a^{p-1} \equiv 1 \ [p]$ \fg{}
devient une égalité. C'est la même chose, dite dans le quotient plutôt que dans
$\mathbb{Z}$, et c'est plus commode : $\mathbb{Z}/p\mathbb{Z}$ est un corps, dont le groupe
multiplicatif a $p - 1$ éléments, et le théorème n'est que le théorème de Lagrange
appliqué à ce groupe.
\end{proof}
\source{https://github.com/Commutator-IO/learn-lean/blob/main/courses/02-lycee/10-specialite-mathematiques-terminale-s/Arithmetique.lean\#L115}{Arithmetique.lean\#L115}

\subsection{Chiffrement affine}

\begin{theoreme}
Chiffrement affine modulo $n$ : si $u$ est inversible modulo $n$, la fonction
$x \mapsto ux + b$ est une bijection de $\mathbb{Z}/n\mathbb{Z}$, et le déchiffrement est
\[ y \mapsto u^{-1}(y - b) . \]
\end{theoreme}

\begin{proof}
Il suffit d'exhiber la réciproque : la fonction $y \mapsto u^{-1}(y - b)$ convient.

En effet, chiffrer puis déchiffrer donne
\[ u^{-1}\bigl((ux + b) - b\bigr) = u^{-1}(ux) = x , \]
et déchiffrer puis chiffrer donne
\[ u\,\bigl(u^{-1}(y - b)\bigr) + b = (y - b) + b = y , \]
les deux simplifications utilisant $u^{-1}u = uu^{-1} = 1$. Une fonction qui admet une
réciproque des deux côtés est bijective.

\medskip\noindent
L'hypothèse d'inversibilité est nécessaire : modulo $26$, chiffrer avec $u = 13$ envoie
toutes les lettres sur deux valeurs seulement, et rien ne se déchiffre.

\medskip\noindent
\textbf{Ce qui n'est pas démontré.} Le chiffrement \textbf{RSA} n'est pas formalisé. Sa
correction \textemdash{} si $ed \equiv 1$ modulo $(p-1)(q-1)$, alors $m^{ed} \equiv m$
modulo $pq$ \textemdash{} s'obtient en appliquant le petit théorème de Fermat modulo $p$
puis modulo $q$, et en recollant par le théorème des restes chinois. Chacune de ces trois
pièces existe dans Mathlib ; c'est leur assemblage, et le traitement des cas où $p$ ou $q$
divise $m$, qui demandent un travail à part entière.
\end{proof}
\source{https://github.com/Commutator-IO/learn-lean/blob/main/courses/02-lycee/10-specialite-mathematiques-terminale-s/Arithmetique.lean\#L123}{Arithmetique.lean\#L123}

\subsection{Énoncés admis}

\noindent
Ce que le programme demande et que ce document ne démontre pas. Les énoncés sont écrits en Lean, pour que le manque se compte ; leur démonstration est admise.

\begin{theoreme}
\emph{Chiffrement RSA.} Si $p$ et $q$ sont deux nombres premiers distincts et si
$ed \equiv 1 \pmod{(p-1)(q-1)}$, alors élever un message à la puissance $e$ puis à la
puissance $d$ le ramène à lui-même modulo $pq$.
\end{theoreme}
\admis
\source{https://github.com/Commutator-IO/learn-lean/blob/main/courses/02-lycee/10-specialite-mathematiques-terminale-s/Arithmetique.lean\#L149}{Arithmetique.lean\#L149}

\section{MatricesEtGraphes}

\noindent
Spécialité de terminale S, section \og{} Matrices et graphes \fg{}. Les matrices sont
celles de Mathlib : le type des tableaux carrés de taille $n$ à coefficients réels,
$A \cdot X$ le produit d'une matrice par un vecteur colonne, $A^{-1}$ l'inverse d'une
matrice de déterminant inversible.

\subsection{Opérations sur les matrices}

\begin{theoreme}
Le produit matriciel est associatif et distributif sur l'addition :
\[ A(BC) = (AB)C , \qquad A(B + C) = AB + AC . \]
\end{theoreme}

\begin{proof}
Les deux règles se lisent coefficient par coefficient. Le coefficient d'indice $(i, j)$ d'un
produit $AB$ est $\sum_k A_{ik} B_{kj}$.

Pour l'associativité, les deux membres donnent la même somme double :
\[ \bigl(A(BC)\bigr)_{ij} = \sum_k A_{ik} \sum_\ell B_{k\ell} C_{\ell j}
   = \sum_{k, \ell} A_{ik} B_{k\ell} C_{\ell j}
   = \sum_\ell \Bigl(\sum_k A_{ik} B_{k\ell}\Bigr) C_{\ell j} = \bigl((AB)C\bigr)_{ij} , \]
l'échange des deux sommes étant licite puisqu'elles sont finies.

Pour la distributivité, c'est celle des réels, sortie du signe somme :
\[ \sum_k A_{ik}\,(B_{kj} + C_{kj}) = \sum_k A_{ik} B_{kj} + \sum_k A_{ik} C_{kj} . \]

C'est ce qui fait des matrices carrées un anneau — non commutatif, l'associativité et la
distributivité ne donnant pas la commutativité.
\end{proof}
\source{https://github.com/Commutator-IO/learn-lean/blob/main/courses/02-lycee/10-specialite-mathematiques-terminale-s/MatricesEtGraphes.lean\#L17}{MatricesEtGraphes.lean\#L17}

\begin{exemple}
Mais il n'est pas commutatif :
\[ \begin{pmatrix} 0 & 1 \\ 0 & 0 \end{pmatrix}
   \begin{pmatrix} 0 & 0 \\ 1 & 0 \end{pmatrix}
   = \begin{pmatrix} 1 & 0 \\ 0 & 0 \end{pmatrix} ,
   \qquad
   \begin{pmatrix} 0 & 0 \\ 1 & 0 \end{pmatrix}
   \begin{pmatrix} 0 & 1 \\ 0 & 0 \end{pmatrix}
   = \begin{pmatrix} 0 & 0 \\ 0 & 1 \end{pmatrix} . \]
\end{exemple}

\begin{explication}
Réfutation par contre-exemple. Si les deux produits étaient égaux, leurs coefficients en
haut à gauche le seraient aussi ; or le premier vaut $1$ et le second $0$. Le calcul des
deux coefficients se fait en développant la somme sur les deux indices, et
$1 \ne 0$ conclut.
\end{explication}
\source{https://github.com/Commutator-IO/learn-lean/blob/main/courses/02-lycee/10-specialite-mathematiques-terminale-s/MatricesEtGraphes.lean\#L22}{MatricesEtGraphes.lean\#L22}

\subsection{Matrice inversible}

\begin{theoreme}
Le déterminant d'une matrice $2 \times 2$ vaut $ad - bc$, et lorsqu'il n'est pas nul la
matrice est inversible :
\[ \begin{pmatrix} a & b \\ c & d \end{pmatrix}^{-1}
   = \frac{1}{ad - bc}\begin{pmatrix} d & -b \\ -c & a \end{pmatrix} . \]
\end{theoreme}

\begin{proof}
Le déterminant d'une matrice $2 \times 2$ est $ad - bc$ par définition.

Pour l'inverse, on vérifie les deux produits, coefficient par coefficient. Par exemple,
celui en haut à gauche du produit de $A$ par
$\frac{1}{ad-bc}\begin{pmatrix} d & -b \\ -c & a\end{pmatrix}$ vaut
\[ \frac{ad - bc}{ad - bc} = 1 , \]
et celui en haut à droite
\[ \frac{-ab + ba}{ad - bc} = 0 . \]
Les quatre coefficients de chacun des deux produits se calculent ainsi, et l'on reconnaît
la matrice identité. Toutes ces simplifications réclament $ad - bc \ne 0$, qui est
l'hypothèse.
\end{proof}
\source{https://github.com/Commutator-IO/learn-lean/blob/main/courses/02-lycee/10-specialite-mathematiques-terminale-s/MatricesEtGraphes.lean\#L33}{MatricesEtGraphes.lean\#L33}

\subsection{Écriture matricielle d'un système}

\begin{theoreme}
Si $A$ est inversible, le système $AX = B$ a une solution et une seule,
\[ X = A^{-1}B . \]
\end{theoreme}

\begin{proof}
\emph{Existence.} En appliquant $A$ au vecteur $A^{-1}B$ :
\[ A\,(A^{-1}B) = (AA^{-1})B = IB = B . \]

\emph{Unicité.} Si $AY = B$, alors en multipliant à gauche par $A^{-1}$ :
\[ A^{-1}B = A^{-1}(AY) = (A^{-1}A)Y = Y . \]

\medskip\noindent
C'est exactement la résolution scolaire, à ceci près que l'inversibilité est écrite comme
celle du \emph{déterminant} : dans la bibliothèque, l'inverse d'une matrice est défini pour
toute matrice, et ne se comporte bien que lorsque le déterminant est inversible. Sans cette
hypothèse, $A^{-1}$ existe encore comme notation, mais $AA^{-1} = I$ est faux : c'est une
valeur de remplissage, pas un inverse.
\end{proof}
\source{https://github.com/Commutator-IO/learn-lean/blob/main/courses/02-lycee/10-specialite-mathematiques-terminale-s/MatricesEtGraphes.lean\#L50}{MatricesEtGraphes.lean\#L50}

\subsection{Puissances d'une matrice}

\begin{theoreme}
Puissances d'une matrice diagonale :
\[ \bigl(\mathrm{diag}(d_1, \dots, d_n)\bigr)^{k}
   = \mathrm{diag}(d_1^{\,k}, \dots, d_n^{\,k}) . \]
\end{theoreme}

\begin{proof}
Le produit de deux matrices diagonales se calcule d'un coup : dans
$\sum_k A_{ik} B_{kj}$, le facteur $A_{ik}$ est nul sauf pour $k = i$, et $B_{kj}$ est nul
sauf pour $k = j$. La somme est donc nulle si $i \ne j$, et vaut $a_i b_i$ sinon :
\[ \mathrm{diag}(a) \times \mathrm{diag}(b) = \mathrm{diag}(a_1b_1, \dots, a_nb_n) . \]

La récurrence sur $k$ suit : $\mathrm{diag}(d)^{k+1} = \mathrm{diag}(d)^{k} \times
\mathrm{diag}(d) = \mathrm{diag}(d^{k}) \times \mathrm{diag}(d) = \mathrm{diag}(d^{k+1})$,
les puissances de la famille $d$ étant prises coefficient par coefficient.

C'est ce qui rend les matrices diagonalisables commodes : élever à la puissance $k$ coûte
$n$ puissances de réels au lieu de $k$ produits de matrices.
\end{proof}
\source{https://github.com/Commutator-IO/learn-lean/blob/main/courses/02-lycee/10-specialite-mathematiques-terminale-s/MatricesEtGraphes.lean\#L62}{MatricesEtGraphes.lean\#L62}

\begin{theoreme}
Puissances par diagonalisation : si $P$ est inversible, alors pour tout $k$,
\[ \bigl(PDP^{-1}\bigr)^{k} = P D^{k} P^{-1} . \]
\end{theoreme}

\begin{proof}
Par récurrence sur $k$.

Au rang $0$, les deux membres valent l'identité : à gauche par convention sur la puissance
nulle, à droite parce que $PP^{-1} = I$, ce qui réclame $P$ inversible.

Supposons la formule vraie au rang $m$. Alors
\begin{align*}
  (PDP^{-1})^{m+1} &= (PDP^{-1})^{m}\,(PDP^{-1}) = P D^{m} P^{-1} \cdot P D P^{-1} \\
   &= P D^{m} (P^{-1}P) D P^{-1} = P D^{m} D P^{-1} = P D^{m+1} P^{-1} ,
\end{align*}
où l'on a simplifié $P^{-1}P = I$ au milieu. C'est tout l'intérêt de la diagonalisation :
les $P^{-1}P$ intermédiaires se télescopent, et il ne reste que les puissances de $D$, qui
se calculent coefficient par coefficient d'après l'énoncé précédent.
\end{proof}
\source{https://github.com/Commutator-IO/learn-lean/blob/main/courses/02-lycee/10-specialite-mathematiques-terminale-s/MatricesEtGraphes.lean\#L68}{MatricesEtGraphes.lean\#L68}

\subsection{Suites de matrices colonnes}

\begin{theoreme}
Suite $U_{n+1} = AU_n + B$ : si $U^{*}$ est un état stable, c'est-à-dire
$U^{*} = AU^{*} + B$, alors
\[ U_n - U^{*} = A^{n}\,(U_0 - U^{*}) . \]
\end{theoreme}

\begin{proof}
Par récurrence sur $n$. Au rang $0$, les deux membres sont $U_0 - U^{*}$, la puissance
nulle de $A$ étant l'identité.

Supposons la formule vraie au rang $m$. En retranchant la relation de récurrence
$U_{m+1} = AU_m + B$ et la relation d'état stable $U^{*} = AU^{*} + B$, les termes $B$
s'éliminent :
\[ U_{m+1} - U^{*} = A U_m - A U^{*} = A\,(U_m - U^{*}) . \]
L'hypothèse de récurrence donne alors
\[ U_{m+1} - U^{*} = A\,\bigl(A^{m}(U_0 - U^{*})\bigr) = A^{m+1}(U_0 - U^{*}) . \]

\medskip\noindent
C'est la méthode du programme : on ramène une suite affine à une suite géométrique en
retranchant l'état stable. La forme explicite $U_n = A^{n}(U_0 - U^{*}) + U^{*}$ s'en
déduit immédiatement.
\end{proof}
\source{https://github.com/Commutator-IO/learn-lean/blob/main/courses/02-lycee/10-specialite-mathematiques-terminale-s/MatricesEtGraphes.lean\#L84}{MatricesEtGraphes.lean\#L84}

\subsection{Graphe probabiliste}

\begin{theoreme}
Graphe probabiliste à deux états, de matrice de transition
$\begin{pmatrix} 1-a & a \\ b & 1-b \end{pmatrix}$ avec $a > 0$, $b > 0$ et $a + b < 2$.
Alors :
\[ \pi = \Bigl(\frac{b}{a+b},\ \frac{a}{a+b}\Bigr) \ \text{est un état stable} , \]
\[ (uM)_0 - \pi_0 = (1 - a - b)\,(u_0 - \pi_0) \quad
   \text{pour tout état } u \text{ avec } u_0 + u_1 = 1 , \]
\[ |1 - a - b| < 1 , \qquad (1-a-b)^{n}\,(u_0 - \pi_0) \xrightarrow[n \to +\infty]{} 0 . \]
\end{theoreme}

\begin{proof}
\emph{L'état $\pi$ est stable.} Sa première coordonnée après une étape vaut
\[ \frac{b}{a+b}(1-a) + \frac{a}{a+b}\,b
   = \frac{b - ab + ab}{a+b} = \frac{b}{a+b} , \]
et la seconde s'en déduit puisque les coordonnées somment à $1$.

\emph{L'écart est contracté.} Pour un état $u$ avec $u_0 + u_1 = 1$, la première
coordonnée après une étape est $u_0(1-a) + u_1 b$, et
\[ u_0(1-a) + u_1 b - \frac{b}{a+b}
   = u_0(1 - a - b) + b - \frac{b}{a+b}
   = (1 - a - b)\Bigl(u_0 - \frac{b}{a+b}\Bigr) , \]
en remplaçant $u_1$ par $1 - u_0$ puis en factorisant. Le facteur $1 - a - b$ est la
seconde valeur propre de la matrice de transition ; la première, égale à $1$, correspond à
l'état stable.

\emph{Le facteur est de module strictement inférieur à $1$.} L'inégalité
$|1 - a - b| < 1$ équivaut à $0 < a + b < 2$, ce que donnent les hypothèses $a > 0$,
$b > 0$ et $a + b < 2$.

\emph{Convergence.} Une puissance $q^{n}$ avec $|q| < 1$ tend vers zéro ; en la multipliant
par la constante $u_0 - \pi_0$, la limite reste nulle. L'écart à l'état stable tend donc
vers zéro, quel que soit l'état initial : c'est la convergence que le programme admet.
\end{proof}
\source{https://github.com/Commutator-IO/learn-lean/blob/main/courses/02-lycee/10-specialite-mathematiques-terminale-s/MatricesEtGraphes.lean\#L104}{MatricesEtGraphes.lean\#L104}

\begin{figure}[h]
\centering
% Construction : deux états A en (0,0) et B en (3.4,0), rayon 0.55.
%   Arcs A→B de poids a, B→A de poids b, boucles A→A de 1−a et B→B de 1−b.
% SVG : x = 60 + 46x, y = 90.
\begin{tikzpicture}[scale=1, >=stealth, every node/.style={font=\small}]
  \node[draw, circle, minimum size=11mm] (A) at (0,0) {$A$};
  \node[draw, circle, minimum size=11mm] (B) at (3.4,0) {$B$};
  \draw[->] (A) to[bend left=28] node[above, font=\footnotesize] {$a$} (B);
  \draw[->] (B) to[bend left=28] node[below, font=\footnotesize] {$b$} (A);
  \draw[->] (A) to[out=160, in=200, looseness=7] node[left, font=\footnotesize] {$1-a$} (A);
  \draw[->] (B) to[out=-20, in=20, looseness=7] node[right, font=\footnotesize] {$1-b$} (B);
\end{tikzpicture}

\caption{Un graphe probabiliste à deux états. Chaque flèche porte la probabilité de la transition, et les deux flèches issues d'un même état ont pour somme $1$.}
\end{figure}

\vfill
\noindent\rule{0.35\linewidth}{0.4pt}\par
\noindent{\footnotesize Leonardo de Moura et Sebastian Ullrich,
\emph{The Lean~4 Theorem Prover and Programming Language},
\emph{Automated Deduction — CADE~28}, Springer, 2021, p.~625--635,
\href{https://doi.org/10.1007/978-3-030-79876-5_37}{doi:10.1007/978-3-030-79876-5\_37}.\par}

\end{document}
